% Post-checkpoint-16 proof packet. Full root/B/C/D proof and primary-source audits passed.
\subsection{Projective representations with vanishing geometric Brauer class}

Throughout this subsection $k$ is algebraically closed of characteristic
$p>0$.  A constant in an $O(r)$ or $O(fr^2)$ estimate depends on the
fixed geometric pair and the fixed bundles explicitly chosen below,
and not on $r$, $f$, a Kummer denominator, or an Artinian parameter ring.
The Frobenius in a one-step construction is always the single
$p$-Frobenius.

\begin{lemma}[Proved: fixed tensor twists and boundary comparisons]
\label{D:fixed-tensor-boundary-comparison}
Let $R$ be a normal two-dimensional local domain and let $P$ be
a fixed rank-one reflexive $R$-module.  Put
\[
 I_P=\operatorname{im}(P\otimes_R P^\vee\longrightarrow R).
                                                        \tag{D.BZ.1}
\]
For every reflexive $M$, this ideal annihilates both the kernel
and the cokernel of
\[
             M\otimes_RP\longrightarrow (M\otimes_RP)^{**}.
                                                        \tag{D.BZ.2}
\]
The same formula holds over an Artinian coefficient base whenever
the source and target extensions have the depth-two Hartogs property.
It does not require reflexive hulls to commute with arbitrary base
change.

In the fixed completed surface charts and the Kummer covers of
Lemma~\ref{D:tame-boundary-summed-index}, allow a one-step
isomorphism on the punctured cover
\[
       F^{[*]}M_j\otimes^{[*]}P_n \simeq M_{j+1},
       \qquad P_n=(P\otimes_RR_n)^{**}.
                                                        \tag{D.BZ.3}
\]
The uniform inverse-matrix pole bound and the summed normalized
boundary-index estimate in that lemma remain valid.  The additional
constant depends on $P$ and the fixed chart only.  In particular,
there is no factor equal to the Kummer denominator or the number
of boundary characters.
\end{lemma}

\begin{proof}
On the puncture, write $z=\sum_\nu p_\nu\ell_\nu\in I_P$,
with $p_\nu\in P$ and $\ell_\nu\in P^\vee$.  If
$N=(M\otimes P)^{**}$, contraction by $\ell_\nu$ gives a map
$N\to M$: the map on the puncture extends by reflexivity.
Consequently
\[
             zh=\sum_\nu p_\nu(\ell_\nu h)
\]
lies in the image of (D.BZ.2).  If $h$ is in its kernel, every
contraction into the torsion-free module $M$ vanishes, and the
same identity gives $zh=0$.  This proves both assertions.
The ideal is the unit ideal at every height-one point, and hence
is either $R$ or is primary for the maximal ideal.

For the relative assertion, use the stated Hartogs equalities
$j_*M_T=M_T$ and $j_*N_T=N_T$ on the punctured relative spectrum.
They extend the contractions.  The displayed identity extends
because the target also has this Hartogs property.  These are
actual maps of relative modules; a base-change assertion about
double duals is not used.  For relative Cohen--Macaulay extensions,
the Hartogs property follows from the Cech complex of a module-regular
sequence of two parameters, and persists along Artinian coefficient
extensions.

Here is a direct denominator comparison for the second assertion.
Use the notation
$A_0=k[[x,a]]\subset R$, $A_n=k[[x,s]]$, $a=s^n$,
and $R_n=R\otimes_{A_0}A_n$ from
Lemma~\ref{D:tame-boundary-summed-index}.
Choose fixed $v\in P$ and $\ell\in P^\vee$ such that
$b=\ell(v)$ is nonzero at the generic point of every
component of $a=0$.  Such a pair exists by prime avoidance:
$P$ is free of rank one at those finitely many points.
Multiplication by $v$ and contraction by $\ell$ extend to
reflexive modules over $R_n$, and both composites are
multiplication by $b$.  Thus multiplication by $v$ has a
cokernel killed by $b$.

Compose the ordinary Frobenius map of $M_j$ with multiplication
by $v$ and (D.BZ.3).  The ordinary-to-reflexive Frobenius
defect of any $R_n$-MCM module is killed by the fixed ideal
of Lemma~\ref{D:Kummer-uniform-Frobenius-defect}; the proof
there is source-module factorization and does not require
an untwisted unit structure.  Relative Frobenius for $R_n/A_n$
contributes the fixed determinant $h(x,s^n)$ used in
(D.TC.5).  Hence, between the free $A_n$-modules of rank
$[R:A_0]r$, the resulting matrix $B_j$ satisfies
\[
 x^K h(x,s^n)b\,\operatorname{coker}(B_j)=0
                                                        \tag{D.BZ.4}
\]
after increasing a fixed $K$.  Replace $b$ by its
$A_0$-norm.  The adjugate of multiplication by $b$ shows
that this norm belongs to $bR$, and its reduction modulo
$a$ is nonzero.  It therefore supplies a fixed scalar
$b_0(x,a)\in A_0$ with $b_0(x,0)\ne0$ in (D.BZ.4).
Modulo $s^{np}$ the inverse of each of
$h(x,s^n)$ and $b_0(x,s^n)$ uses at most $p$ terms
of its geometric expansion in $s^n$.  Thus
\[
       B_j^{-1}\bmod s^{np}
            \in x^{-L}\operatorname{Mat}
                         (A_n/(s^{np}))
                                                        \tag{D.BZ.5}
\]
for one fixed $L$.

The element $v$ has character zero.  After choosing homogeneous
$A_n$-bases, the invariant Frobenius comparison is consequently
the same diagonal map with entries $a^{\lfloor pj/n\rfloor}$
as before.  The resulting diagonal quotient maps to the pure
coarse cokernel.  It is an isomorphism after inverting $x$:
there $P$ is trivialized by $v$ modulo every required power of
$a$.  Torsion-freeness over $k[[x]]$ makes this an inclusion.
Equation (D.BZ.5) bounds the reverse inclusion, exactly as in
(D.TC.4).  Intersecting one generic filtration with both
lattices bounds the \emph{sum} of all grade indices by their
total index.  Passing to the fixed normalization of $R/(a)$
multiplies that sum by its fixed rank over $k[[x]]$.

Finally, changing $v$ or a local line trivialization changes
the normalized character lattices by the divisors of fixed
rational functions.  These produce fixed comparison divisors
on the normalized boundary curves.  This accounts for their
global gluing.  We have not asserted that $P_n/sP_n$ is a
base-change-compatible reflexive hull.
\end{proof}

\begin{proposition}[Proved: cyclic tensor twists in the quadratic tame atlas]
\label{D:finite-tensor-cyclic-quadratic-atlas}
Fix the normal proper surface pair $(S,T)$ of
Theorem~\ref{D:tame-global-quadratic-representation-bound}.
Enlarge a fixed finite set $\Sigma$ so that
$S^\circ=S\setminus\Sigma$ is smooth and the reduced
original divisor $D$ is smooth with disjoint components there.
Put $U^\circ=S^\circ\setminus D$.
Fix a smooth projective resolution $\pi:Y\to S$, an
isomorphism over $S^\circ$, the exceptional cycle $Z$,
and the strict boundary curves $C_i$ used in that theorem.
Let $\mathscr L$ be a finite collection of line bundles on $Y$.

Allow the coefficient stages to satisfy, on the tame root
stack over $(S^\circ,D)$,
\[
       F^*\mathcal E_j\otimes L_j
                  \simeq\mathcal E_{j+1},
       \qquad L_j\in\mathscr L,\quad j\in\mathbf Z/f\mathbf Z.
                                                        \tag{D.BZ.6}
\]
Here the boundary weights lie in $[0,1)$ with denominators
prime to $p$, and the corresponding ordinary lattices are
extended reflexively to $S$.  There are constants, independent
of the length and weights of the cycle, such that:
\begin{enumerate}
\item minimal extensions $F_j$ on $Y$ have fixed exceptional
and saturated boundary-grade slope windows; $c_1(F_j)$ has
norm $O(r)$ in $N^1(Y)_{\mathbf R}$; and
$|\chi(Y,F_j)|=O(r)$;
\item each fixed twist, one Frobenius pullback, and each of the
fixed auxiliary sheaves in the filtered-diagram construction
has all its cohomology dimensions bounded by $O(r)$;
\item a word of $f$ tensor and boundary-digit types has
polynomial witnesses in $O(fr^2)$ variables and equations of
bounded degree.  Valid unit-cycle isomorphism classes are
open and closed in its witness locus, and their number is at
most $\exp(Cfr^2)$ after summing over the multiplicity arrays.
\end{enumerate}
The constants may depend on $\mathscr L$.  No bound for
an arbitrarily varying tensor line bundle is asserted.
\end{proposition}

\begin{proof}
Write $P_L=(\pi_*L)^{**}$ for the corresponding fixed
rank-one reflexive sheaf on $S$.  There are only finitely
many such sheaves.  On $S^\circ$ the one-step coarse
map in (D.BZ.6) extends reflexively as
\[
 G_j=\bigl(F_S^{[*]}E_j\otimes P_{L_j}\bigr)^{**}
                   \longrightarrow E_{j+1}.
                                                        \tag{D.BZ.7}
\]
It is injective and has a pure one-dimensional cokernel $T_j$.
Its generic exponents on every original divisor lie between
$0$ and $p-1$.  Purity follows from the depth lemma, since
both terms are reflexive on a normal surface.

\emph{Fixed poles, slopes, and first Chern classes.}
Minimal extensions have the fixed exceptional windows of
Langer's extension lemma
\citep[Definition~1.2 and Lemma~1.3]{C_Langer2026}.
One Frobenius pullback has fixed enlarged windows, by the
curve Frobenius bounds already used in
Proposition~\ref{D:global-periodic-reflexive-boundedness}.
Tensoring by a member of $\mathscr L$ shifts these windows
by one of finitely many fixed degrees.  The bounded-pole
argument therefore gives one $B$ for both maps
\[
 \begin{split}
 F_Y^*F_j\otimes L_j&\longrightarrow F_{j+1}(BZ),\\
 F_{j+1}&\longrightarrow
 F_Y^*F_j\otimes L_j\bigl((p-1)\textstyle\sum_iC_i+BZ\bigr).
 \end{split}                                            \tag{D.BZ.8}
\]
These maps are the genuine maps off the exceptional set;
the second is the ordinary inverse with the fixed allowed
boundary pole.

Taking maximal slopes, and then dualizing the second map,
gives
\[
 p\mu_{\max}(F_j)\le\mu_{\max}(F_{j+1})+C,\qquad
 p\mu_{\max}(F_j^\vee)\le\mu_{\max}(F_{j+1}^\vee)+C.
                                                        \tag{D.BZ.9}
\]
The same inequalities hold on a fixed smooth very ample
curve not contained in the determinant-zero divisors.
The maps restricted to that curve are generically injective
and hence injective.  Taking a maximum around the finite
cycle bounds all four kinds of slopes uniformly.

The determinant of (D.BZ.8) gives
\[
 c_1(F_{j+1})-p c_1(F_j)
                 =r c_1(L_j)+D_j,                       \tag{D.BZ.10}
\]
where $D_j$ is supported on the fixed $C_i$ and the
exceptional curves, and each coefficient has absolute
value at most $Cr$.  This follows from the two pole bounds
in (D.BZ.8); at an original boundary component it is also
the sum of the $r$ exponents in $[0,p-1]$.
After one complete cycle, (D.BZ.10) expresses $c_1(F_j)$
as a geometric weighted sum divided by $p^f-1$.
All its summands belong to a fixed finite-dimensional
span and have norm $O(r)$.  Hence $\|c_1(F_j)\|=O(r)$,
uniformly in $f$.  This argument takes place on $Y$ and
does not assume that $P_L$ is Cartier on $S$.

Normalizing (D.BZ.8) on the boundary root charts gives
exactly the two maps of
Lemma~\ref{D:saturated-boundary-two-sided-slopes}, with
the additional line twist $L_j|_{C_i}$.  At $C_i\cap Z$
this line bundle is locally trivial.  The intersection
of the localizations in its two distinct parameters
therefore proves that normalization preserves the same
exceptional pole bound, for the map and its inverse.
On a character grade the two inequalities become
\[
 p\mu_{\max}(G_{j,i,a})
       \le\mu_{\max}(G_{j+1,i,pa})+C,\qquad
 p\mu_{\max}(G_{j,i,a}^{\vee})
       \le\mu_{\max}(G_{j+1,i,pa}^{\vee})+C.
                                                        \tag{D.BZ.11}
\]
The fixed constant absorbs $\deg(L_j|_{C_i})$ and the
normal digit twist.  The pairs consisting of stage and
character form finite cycles.  Thus every saturated
grade of the full flag inside the \emph{same} $F_j|_{C_i}$
has fixed upper and lower slope bounds.

\emph{The Euler estimate, including the fixed tensor defect.}
Lemma~\ref{D:fixed-tensor-boundary-comparison} extends the
local determinant, inverse, and summed-index calculation
to (D.BZ.7).  On normalized complete boundary curves the
forward character maps have only the additional degrees
of the finitely many $P_L$ trivializations.  Taking maxima
around the stage-character cycles gives the same uniform
upper-slope estimates as (D.TC.10).  Each character is
repeated at most $p-1$ times.  The one-filtration index
calculation bounds all positive corrections together by
$Cr$.  The proof of Lemma~\ref{D:tame-canonical-Euler}
therefore gives
\[
                         \chi(S,T_j)\le Cr.              \tag{D.BZ.12}
\]

We also justify quantitatively the tensor change in the
source of (D.BZ.7).  The relative Euler calculation in the
proof of Proposition~\ref{D:global-periodic-reflexive-boundedness}
uses only fixed exceptional slope windows.  For any bundle
$A$ with those windows and
$E=(\pi_*A)^{**}$, it gives
\[
                         |\chi(Y,A)-\chi(S,E)|\le Cr.
                                                        \tag{D.BZ.13}
\]
Indeed $E=\pi_*A(BZ)$ for a fixed $B$; the quotient is
bounded by the cohomology of one fixed exceptional
thickening.  Formal functions and the same layer
vanishing bound $R^1\pi_*A$ by another fixed thickening.
Their component cohomology is $O(r)$ from the upper and
lower slope windows.  This assertion is not restricted
to periodic $E$.

Apply (D.BZ.13) first to $A=F_Y^*F_j$ and then to
$A\otimes L_j$.  Their reflexive pushforwards are
$F_S^{[*]}E_j$ and $G_j$, respectively.  Smooth-surface
Riemann--Roch and (D.BZ.10) give
\[
 \begin{split}
 \chi(S,G_j)-\chi(S,F_S^{[*]}E_j)
 &=\chi(Y,A\otimes L_j)-\chi(Y,A)+O(r)\\
 &=c_1(A)\cdot c_1(L_j)
   +\frac r2\bigl(c_1(L_j)^2-K_Yc_1(L_j)\bigr)+O(r)
   =O(r).
 \end{split}                                            \tag{D.BZ.14}
\]
No formula asserting exact multiplicativity of a
non-Cartier reflexive tensor product on $S$ is used.
Together with (D.BZ.12), normal-surface Riemann--Roch
and Frobenius scaling
\citep[Theorem~0.1]{C_Langer2026}
\citep[Theorem~4.5(2)]{C_Langer2025}
now yield
\[
       \operatorname{ch}_2(E_{j+1})
                   -p^2\operatorname{ch}_2(E_j)\le Cr.
                                                        \tag{D.BZ.15}
\]
The linear first-Chern terms are $O(r)$ by (D.BZ.10).
Cycling the inequality gives
$\operatorname{ch}_2(E_j)\ge-Cr$, with the sign shown.

The slope bounds (D.BZ.9) give $h^0(S,E_j)\le Cr$ by
the fixed-pole comparison and restriction to the fixed
very ample curve.  Proper duality bounds $h^2(S,E_j)$
by $Cr$: $\operatorname{Hom}(E_j,\omega_S)$ injects
into sections of $F_j^\vee$ tensored by a fixed line
bundle and a fixed exceptional pole, and the dual
slope bound applies.  The lower Riemann--Roch estimate
from (D.BZ.15) consequently gives
$h^1(S,E_j)\le Cr$ and $|\chi(S,E_j)|\le Cr$.
Equation (D.BZ.13) proves $|\chi(Y,F_j)|\le Cr$.
The same fixed-curve restrictions, duality, and smooth
Riemann--Roch show
$h^i(Y,F_j(t))\le C_t r$ for all fixed $t$ and $0\le i\le2$.
They also bound one Frobenius pullback and each fixed
tensor auxiliary: their first Chern characters and
second Chern characters are $O(r)$ by (D.BZ.10), the
Euler estimate, and one-step Frobenius scaling.
The threshold-transform auxiliary sheaves have the
same bounds, because there are at most $p$ transforms
and all full-flag grades have the fixed windows
in (D.BZ.11).

\emph{Polynomial witnesses and Artinian effectiveness.}
Use the monads, filtered linear curve resolutions,
and fixed-column derived Hom construction of
Theorem~\ref{D:tame-global-quadratic-representation-bound}.
Insert tensoring by $L_j$ at its one-step Frobenius
node.  This is one of finitely many fixed operations:
tensor the fixed bundle terms of its complexes by
$L_j$, and use the same coherent tensor diagrams
to identify it with each pole-twisted auxiliary.
Their total multiplicities are $O(r)$ by the preceding
cohomology bounds.  Their lengths and bundle types
are fixed.  Therefore the number of scalar variables
and equations per stage remains $O(r^2)$ and their
degrees remain bounded.

The boundary digit construction is unchanged:
digits belong to $[0,p-1]$, and the permutation is
increasing inside each equal-digit block.  Tensoring
by a line has no effect on its canonical block
splitting or on the ceiling-valuation matrix
criterion.  A fixed tensor type word adds at most
$|\mathscr L|^f$ choices.  The boundary words and
multiplicity arrays have the same $\exp(O(fr^2))$
bound as before.

For completeness, the open-class argument is valid
with specified Artinian points.  Composing the cycle
gives
\[
       F^{f*}\mathcal E_0\otimes M\simeq\mathcal E_0,
       \qquad M=\bigotimes_{j=0}^{f-1}
                         L_j^{\,p^{f-1-j}}.
                                                        \tag{D.BZ.16}
\]
On the $(p^f-1)$st root gerbe of $M$, with universal
line $\mathcal L$, this is a unit structure on
$\mathcal E_0\otimes\mathcal L$.
The gerbe is used over the smooth open and its
boundary root stack; it is not extended as an
Azumaya algebra across the singular points of $S$.
Finite etale fixed-frame torsors are invariant
under Artinian nilpotent thickenings.  A specified
unit isomorphism therefore lifts uniquely there,
as is checked on an etale atlas and its descent
diagram.  Untwisting gives the specified maps of
all ordinary stages and their boundary flags.

The fixed exceptional-window pole bound extends
these maps and their inverses to $Y$ with one $B$,
independent of Artinian length.  The layer-vanishing
proof works over an Artinian base by its
maximal-ideal filtration.  Flag preservation
extends at $C_i\cap Z$ because the relative
quotients are vector bundles on $C_i$, so its
local parameter remains a non-zero-divisor.
The actual fixed line twists commute with base
change.  If a reflexive description downstairs
is used, the first part of
Lemma~\ref{D:fixed-tensor-boundary-comparison}
gives the needed tensor comparison through
Hartogs, rather than through an assumed
base-change theorem for reflexive hulls.

Proper bounded-pole Hom schemes with inverse and
cycle equations thus represent a unit-Isom functor
of finite presentation which lifts every specified
small Artinian point.  Its image is open by the
infinitesimal smoothness criterion, as in the
preceding quadratic theorem.  Distinct unit-cycle
classes give disjoint open fibres covering the
witness scheme.  They form a finite clopen
partition, with at most as many members as
irreducible components.  The witness scheme is
an open subset of an affine scheme defined by
$O(fr^2)$ bounded-degree equations.  Bezout and
the array/word estimates prove the last assertion.
\end{proof}

\begin{theorem}[Proved: quadratic growth for Brauer-zero projective representations]
\label{D:Brauer-zero-projective-representation-growth}
Let $S$ be a normal proper connected surface over $k$,
let $T\subsetneq S$ be closed, and put $Q=Q(S,T)$.
For a representation
$\rho:Q\to\operatorname{PGL}_r(\mathbf F_q)$,
$q=p^f$, denote by $\operatorname{br}(\rho)$ the
class in $H^2(U,\mathbf G_m)$ of its induced
algebraic $\operatorname{PGL}_r$-torsor on $U=S-T$.
There is a constant $C_{S,T}$ such that
\[
 \#\left(
   \{\rho:Q\to\operatorname{PGL}_r(\mathbf F_q)
                   \mid \operatorname{br}(\rho)=0\}
                   /\operatorname{PGL}_r(\mathbf F_q)
    \right)
               \le q^{C_{S,T}r^2}                       \tag{D.BZ.17}
\]
for every $r,f\ge1$.  The scalar centre is not required
to have bounded order.
\end{theorem}

\begin{proof}
Enlarge the fixed finite omitted set to $\Sigma$
as in the preceding proposition.  The resulting
original-boundary tame group $Q^\circ$ surjects
onto $Q$.  Conjugacy classes inject under this
surjection, and a zero Brauer class stays zero
on $U^\circ$.  It suffices to count for the
smooth $U^\circ$.

\emph{A fixed finite set of scalar digit twists.}
Restriction $\operatorname{Pic}(Y)\to
\operatorname{Pic}(U^\circ)$ is surjective:
extend a Cartier divisor on the smooth open
by its closure, which is Cartier on the
regular surface $Y$.  The Neron--Severi
group of the smooth projective $Y$ is
finitely generated
\citep[Expose XIII, Theorem~5.1 and proof,
printed pp.~650--651]{SGA6BrauerZeroD}.
The group of $k$-points of $\operatorname{Pic}^0(Y)$
is divisible by every integer prime to $p$:
its reduced connected group is an abelian
variety and multiplication by such an integer
is surjective
\citep[Proposition~39.9.11(3)--(4), Tag 03RP;
Lemma~39.9.8 and proof, Tag 0BFG]{StacksBrauerZeroD}.
Choose line bundles
$L_1,\ldots,L_b$ representing generators
of the Neron--Severi group, including its
finite torsion generators.  For every
$n$ prime to $p$, every class in
$\operatorname{Pic}(U^\circ)/n$ has a
representative
\[
          M_{\mathbf a}=\bigotimes_{i=1}^b L_i^{a_i}
                         \bigm|_{U^\circ},
               \qquad 0\le a_i<n.                       \tag{D.BZ.18}
\]
Only the smooth projective Picard group is used
here; no identical assertion about a complete
local Picard group is being assumed.

Put $n=q-1$.  The finite central extension
\[
 1\longrightarrow\mathbf F_q^\times
  \longrightarrow\operatorname{GL}_r(\mathbf F_q)
  \longrightarrow\operatorname{PGL}_r(\mathbf F_q)
  \longrightarrow1
                                                        \tag{D.BZ.19}
\]
gives the lifting gerbe of $\rho$.
Identify $\mathbf F_q^\times$ with $\mu_n(k)$.
Its class maps to $\operatorname{br}(\rho)$.
The Kummer sequence
\citep[Lemma~59.28.1, Tag 03PK]{StacksBrauerZeroD}
therefore makes the
lifting gerbe, when this class is zero,
bandedly isomorphic to the root gerbe
$\sqrt[n]{M_{\mathbf a}}$ for some
(D.BZ.18).  The tautological lift is a
finite $\mathbf F_q$ local system of
scalar weight one on this gerbe.
The root-gerbe untwisting equivalence
\citep[Sections~2.2.2 and 3.1.1;
Lemmas~3.1.1.7--3.1.1.8 and
3.1.1.11--3.1.1.12]{Lieblich2008BrauerZeroD}
writes its vector bundle as
$E_0\otimes\mathcal L$, where
$\mathcal L^n=M_{\mathbf a}$.
Its $q$-unit structure is precisely
\[
                 F^{f*}E_0\otimes M_{\mathbf a}
                                  \simeq E_0.           \tag{D.BZ.20}
\]

The root gerbe of $M_{\mathbf a}$ extends
over $S^\circ$, including its original
boundary DVRs.  The projective representation
is tame there.  Its wild pro-$p$ inertia
can therefore act on the lift only by
the prime-to-$p$ scalar band, and must act
trivially.  Hence the lifted unit bundle
has a tame parabolic extension on the
boundary root stack.  Conversely a
weight-one unit bundle with this extension
projectivizes to a representation of
$Q^\circ$.  No tameness is imposed at
the additional exceptional valuations.

\emph{Replacing the growing scalar root by fixed digits.}
For each coordinate of $\mathbf a$ define
the cyclic residues and digits by
\[
 a_{i,j+1}=p a_{i,j}-n c_{i,j},\qquad
 0\le a_{i,j}<n,\quad 0\le c_{i,j}\le p-1.
                                                        \tag{D.BZ.21}
\]
Since $p^f\equiv1\pmod n$, the residues
are cyclic.  Let
$L_j=\bigotimes_i L_i^{c_{i,j}}$.
These bundles belong to the one fixed
finite collection
$\mathscr L=\{\bigotimes_iL_i^{c_i}:0\le c_i<p\}$.
Writing
$b_{i,j}=\lfloor p^j a_{i,0}/n\rfloor$
and
$E_j=F^{j*}E_0\otimes\bigotimes_iL_i^{b_{i,j}}$
on the smooth open gives
\[
              F^*E_j\otimes L_j\simeq E_{j+1},
                         \qquad E_f=E_0.               \tag{D.BZ.22}
\]
Their canonical parabolic lattices extend
reflexively to $S$.  Thus every representation
being counted supplies data of the preceding
proposition.

Conversely a word of the fixed scalar digits
determines
$a_{i,0}=\sum_{j=0}^{f-1}p^{f-1-j}c_{i,j}$.
Discard the all-$(p-1)$ word in any coordinate
if it gives $a_{i,0}=n$; no chosen representative
in (D.BZ.18) is lost.  A valid scalar word
and valid boundary word reconstruct the
root gerbe and a weight-one unit bundle
from (D.BZ.22), by (D.BZ.16).
Its projectivization gives the desired
finite-field projective representation.
Choices of the root-gerbe identification
or intermediate edge isomorphisms can only
increase the count of unit-cycle classes.
There are at most $p^{bf}=q^b$ scalar words,
and no equation of degree $q$ or $n$ has
been inserted into the one-step atlas.

For clarity, the finite-class comparison
has no hidden injectivity assertion for
the band of an arbitrary gerbe.  A
fundamental group of a root gerbe may
have a kernel which is only an image of
$\mu_n$.  If a weight-one finite unit
bundle exists, its projective torsor
$P$ identifies that gerbe with the
gerbe of lifts of $P$ through (D.BZ.19).
For the finite monodromy image $\Gamma$, now
let $P$ denote the connected $\Gamma$-torsor
reduction determined by that image, rather
than the full projective-frame torsor.
For its inverse image
$\widetilde\Gamma$ in
$\operatorname{GL}_r(\mathbf F_q)$,
the gerbe is the quotient
$[P/\widetilde\Gamma]$ and its tame
fundamental group is
\[
                  \widetilde\Gamma
                          \times_\Gamma Q^\circ.
                                                        \tag{D.BZ.23}
\]
Indeed its finite covers are exactly
finite covers of $P$ with
$\widetilde\Gamma$-equivariant descent;
the same descent identifies the displayed
pullback group.  Wild inertia lifts
uniquely through the prime-to-$p$
central kernel, so quotienting by its
normal closure gives exactly (D.BZ.23).

A finite central extension of a small
profinite group is small.  For a finite
target, fix its finite image $H$, the
image of the central kernel, and the
restriction to that kernel.  There
are finitely many induced maps from
the quotient to $H$ modulo that central
image.  Two lifts with the same data
differ by a character of the quotient
to that finite central image, again
with only finitely many possibilities.
Theorem~\ref{D:fixed-boundary-small}
therefore supplies the appropriate
finite-class interpretation whenever
the root gerbe has such an object;
empty root classes contribute nothing.
The quantitative open-fibre argument
of the preceding proposition is
stronger and already gives the needed
component count.

For every scalar and boundary word that
argument bounds the unit-cycle classes
by $\exp(Cfr^2)$, including all monad
and flag arrays.  Multiplying by
$q^b$ and the already bounded boundary
type choices only changes $C$.
Projectivization and the initial
surjection to $Q$ prove (D.BZ.17).
\end{proof}

\begin{remark}[Exact scope of the reduction]
This controls projective representations whose geometric
Brauer class vanishes, without a bound on their scalar
centre.  Their obstruction in $H^2(U^\circ,\mu_{q-1})$
may be a nonzero Kummer class; an ordinary finite-field
linear lift is not assumed.  The earlier bounded-centre
argument killed finitely many coefficient orders on one
fixed open subgroup.  Finiteness of $H^2$ separately at
each order does not supply that subgroup uniformly as
$q-1$ grows.  The finite-digit construction supplies this
additional uniformity for the Kummer classes.
It does not control the nonzero Brauer classes,
construct an Azumaya extension across the boundary, or
prove finite generation of the full fundamental group.
Complete local punctures are not included by the
projective Picard argument above.
\end{remark}
